\chapter{Training Deep Networks}

\section{Initialization Strategies}

Training deep neural networks involves optimizing highly nonconvex objectives in high-dimensional parameter spaces. The choice of initialization plays a crucial role in determining whether optimization succeeds.

\medskip

\noindent
\textbf{Random initialization.}  
Parameters are typically initialized randomly:
\[
W_\ell^{(0)} \sim \text{distribution}, \quad b_\ell^{(0)} = 0 \ \text{or small values}.
\]

\medskip

\noindent
\textbf{Why not initialize to zero?}
\begin{itemize}
    \item Symmetry: all units would behave identically
    \item No meaningful learning would occur
\end{itemize}

\medskip

\noindent
\textbf{Variance preservation.}  
A key principle is to choose initialization so that activations neither vanish nor explode as they propagate through layers.

\medskip

\noindent
\textbf{Example (Xavier/Glorot initialization).}
\[
\mathrm{Var}(W_{ij}) \approx \frac{2}{n_{\text{in}} + n_{\text{out}}}.
\]

\medskip

\noindent
\textbf{Example (He initialization).}
\[
\mathrm{Var}(W_{ij}) \approx \frac{2}{n_{\text{in}}}.
\]

\medskip

\noindent
\textbf{Insight.}  
Proper initialization maintains stable signal propagation across layers.

\medskip

\noindent
\textbf{Effect on training.}
\begin{itemize}
    \item Good initialization $\rightarrow$ faster convergence
    \item Poor initialization $\rightarrow$ slow or failed training
\end{itemize}

\section{Vanishing and Exploding Gradients}

A major difficulty in training deep networks arises from the behavior of gradients.

\medskip

\noindent
\textbf{Gradient propagation.}  
Backpropagation involves repeated multiplication by weight matrices and derivatives:
\[
\frac{\partial L}{\partial z^{(\ell)}} = \frac{\partial L}{\partial z^{(\ell+1)}} W_{\ell+1}^\top \cdot \sigma'(z^{(\ell)}).
\]

\medskip

\noindent
\textbf{Vanishing gradients.}
\begin{itemize}
    \item Gradients shrink exponentially with depth
    \item Early layers learn very slowly
\end{itemize}

\medskip

\noindent
\textbf{Exploding gradients.}
\begin{itemize}
    \item Gradients grow exponentially
    \item Leads to instability and numerical issues
\end{itemize}

\medskip

\noindent
\textbf{Causes.}
\begin{itemize}
    \item Repeated multiplication of small or large values
    \item Saturating activation functions (e.g., sigmoid)
\end{itemize}

\medskip

\noindent
\textbf{Mitigation strategies.}
\begin{itemize}
    \item Careful initialization
    \item Use of non-saturating activations (e.g., ReLU)
    \item Normalization methods
    \item Gradient clipping
\end{itemize}

\medskip

\noindent
\textbf{Insight.}  
Stable gradient flow is essential for effective learning in deep networks.

\section{Optimization Challenges}

Deep networks introduce optimization challenges beyond those encountered in classical models.

\medskip

\noindent
\textbf{Nonconvexity.}  
The loss function has many local minima and saddle points.

\medskip

\noindent
\textbf{High dimensionality.}  
The parameter space can contain millions or billions of variables.

\medskip

\noindent
\textbf{Saddle points.}
\begin{itemize}
    \item More common than poor local minima
    \item Can slow down optimization
\end{itemize}

\medskip

\noindent
\textbf{Flat vs sharp regions.}
\begin{itemize}
    \item Flat regions: slow progress
    \item Sharp regions: instability
\end{itemize}

\medskip

\noindent
\textbf{Ill-conditioning.}  
Different directions may have very different curvature, making optimization difficult.

\medskip

\noindent
\textbf{Practical strategies.}
\begin{itemize}
    \item Use stochastic optimization (SGD)
    \item Apply momentum or adaptive methods
    \item Normalize inputs and activations
\end{itemize}

\medskip

\noindent
\textbf{Insight.}  
Despite theoretical complexity, simple optimization methods often work well in practice due to structure in real-world problems.

\section{Empirical Training Behavior}

Training deep networks often exhibits characteristic patterns that are not fully captured by theory.

\medskip

\noindent
\textbf{Training dynamics.}
\begin{itemize}
    \item Rapid initial decrease in loss
    \item Slower convergence in later stages
\end{itemize}

\medskip

\noindent
\textbf{Generalization behavior.}
\begin{itemize}
    \item Training loss may approach zero
    \item Test performance may still be good
\end{itemize}

\medskip

\noindent
\textbf{Double descent (informal).}  
Increasing model size can first increase and then decrease test error.

\medskip

\noindent
\textbf{Role of overparameterization.}
\begin{itemize}
    \item Large models are easier to optimize
    \item Many solutions achieve low training error
\end{itemize}

\medskip

\noindent
\textbf{Implicit regularization.}  
Optimization dynamics favor certain solutions, often with good generalization properties.

\medskip

\noindent
\textbf{Learning curves.}
\begin{itemize}
    \item Plot training and validation error over time
    \item Used to diagnose underfitting and overfitting
\end{itemize}

\medskip

\noindent
\textbf{Practical observations.}
\begin{itemize}
    \item Larger datasets improve performance
    \item Deeper models can capture more structure
    \item Careful tuning is often required
\end{itemize}

\medskip

\noindent
\textbf{Insight.}  
Empirical behavior provides important guidance when theory is incomplete.

\section{A Unifying Perspective}

Training deep networks involves several interacting components:

\begin{itemize}
    \item \textbf{Initialization:} sets the starting point
    \item \textbf{Gradient flow:} determines signal propagation
    \item \textbf{Optimization:} navigates the loss landscape
    \item \textbf{Data and model size:} influence generalization
\end{itemize}

\medskip

\noindent
\textbf{Core challenge.}  
Ensuring stable and efficient training in high-dimensional, nonconvex settings.

\medskip

\noindent
\textbf{Key insight.}  
Successful training relies on aligning architecture, initialization, and optimization.

\section{Outlook}

This chapter examined the practical challenges of training deep networks:
\begin{itemize}
    \item initialization strategies,
    \item vanishing and exploding gradients,
    \item optimization difficulties,
    \item empirical training behavior.
\end{itemize}

\medskip

In the next chapter, we study activation functions and architectural design choices, which further influence the behavior of deep networks.

\medskip

\noindent
\textbf{Guiding principle.}  
Training deep networks requires balancing mathematical principles with empirical insight.